%==============================================================================
%  APPENDIX - Operations and service commitments
%  SLOs, backups, reversibility: what makes the system deliverable
%==============================================================================

%==============================================================================
\chapter{Appendix: Operations and service commitments}
\label{chap:exploitation}
%==============================================================================
A system delivered without a measurable service commitment cannot be operated:
with no objective, no degradation can be qualified as an incident. This appendix
sets the service objectives, the backup policy and the conditions of
reversibility. It is the reference for the acceptance criteria of
chapter~\ref{chap:delais}.

%------------------------------------------------------------------------------
\section{Service level objectives (SLO)}
%------------------------------------------------------------------------------
The objectives are sized for a \textbf{single-server} deployment, in line with
the operations decision adopted. They are deliberately realistic: promising
99.99~\% without a redundant architecture would be an unkeepable commitment.

\begin{xltabular}{\textwidth}{>{\bfseries}p{4.2cm} p{3.4cm} X}
\toprule
\textbf{Indicator} & \textbf{Objective} & \textbf{Detail} \\
\midrule
\endhead
Availability & 99.5~\% monthly & About 3~h~40 of tolerated downtime per month, excluding planned maintenance. \\
\addlinespace
Planned maintenance & 4~h/month maximum & Announced 7 days in advance, outside the production season where possible. \\
\addlinespace
Latency (p95) & $<$ 500~ms & On consultation screens, at nominal load. \\
\addlinespace
Latency (p99) & $<$ 1,500~ms & Tolerates heavy analytical queries. \\
\addlinespace
Error rate & $<$ 1~\% & 5xx responses over total requests. \\
\addlinespace
\textbf{RTO} (recovery time) & $<$ 4~h & Maximum delay to restore service after a major incident. \\
\addlinespace
\textbf{RPO} (data loss) & $<$ 1~h & Maximum acceptable loss, guaranteed by the backup frequency. \\
\bottomrule
\end{xltabular}

\begin{encadre}
\textbf{Reference nominal load:} 200 concurrent users at seasonal peak, 10,000
instrumented hives, about 1 million measurements ingested per day. The latency
objectives hold only within these limits; beyond them, the sizing must be
revised.
\end{encadre}

%------------------------------------------------------------------------------
\section{Backup and restoration}
%------------------------------------------------------------------------------
\subsection{Policy}
\begin{itemize}
  \item \textbf{Continuous backup} of PostgreSQL transaction logs (WAL
        archiving), guaranteeing the one-hour RPO.
  \item \textbf{Daily full backup}, encrypted at rest.
  \item \textbf{Off-site storage} on a medium separate from the production
        server --- a backup stored on the machine it protects protects nothing.
  \item \textbf{Retention}: 7 daily backups, 4 weekly, 12 monthly.
  \item \textbf{Visit photographs} and the \texttt{ConfigZumm.ini} file are
        within the backup scope, on the same footing as the database.
\end{itemize}

\subsection{Restoration test}
\begin{encadre}
\textbf{A backup never restored is not a backup.} Full restoration is tested
\textbf{monthly} on a dedicated environment, and the report is archived. A
successful restoration test is a blocking acceptance criterion.
\end{encadre}

The test verifies three points: the integrity of the restored data, compliance
with the announced RTO, and the ability of an operator to carry out the
procedure using the runbook alone.

%------------------------------------------------------------------------------
\section{Monitoring and alerting}
%------------------------------------------------------------------------------
\begin{xltabular}{\textwidth}{>{\bfseries}p{4.2cm} p{3.4cm} X}
\toprule
\textbf{Signal} & \textbf{Alert threshold} & \textbf{Expected action} \\
\midrule
\endhead
Application unavailable & 2 consecutive failures & Immediate intervention \\
Disk space & $>$ 80~\% & Extension or purge within 48~h \\
Latency p95 & $>$ 500~ms over 15~min & Analyse the offending query \\
5xx error rate & $>$ 1~\% over 5~min & Immediate intervention \\
Backup failure & 1 occurrence & Same-day intervention \\
TLS certificate & expiry $<$ 30~d & Renewal \\
Sensor ingestion lag & $>$ 1~h & Check the MQTT broker \\
\bottomrule
\end{xltabular}

Monitoring relies on \textbf{OpenTelemetry}, Prometheus and Grafana, in line
with the stack of appendix~\ref{chap:technologies}.

%------------------------------------------------------------------------------
\section{Reversibility}
%------------------------------------------------------------------------------
The client must be able to take the system back without depending on the
supplier. The following are handed over at delivery, and on any later request:

\begin{itemize}
  \item the \textbf{full data export} in an open, documented format (SQL and
        CSV), including photographs and attachments;
  \item the \textbf{source code} and its version history;
  \item the \textbf{operations documentation} (runbook): start, stop,
        restoration, certificate rotation, reading the alerts;
  \item the \textbf{database schemas} and migration scripts;
  \item the \textbf{redeployment procedure} from a clean repository, verified
        during acceptance.
\end{itemize}

\begin{encadre}
\textbf{Verification:} reversibility is demonstrated, not declared. A clean
redeployment from the repository, onto a fresh machine, is performed during
acceptance.
\end{encadre}

%------------------------------------------------------------------------------
\section{Knowledge transfer}
%------------------------------------------------------------------------------
Two days of transfer are planned at delivery, covering routine administration,
reading the monitoring, the restoration procedure and first-level
troubleshooting. The runbook is the supporting material and remains the
reference after the transfer.
